\chapter{项目创新点}

本项目的创新不建立在动态障碍预测、不确定性建模、运动基元、通信质量评价或协同模式切换等单项技术之上，而是针对现有方法在具体应用环节中的不足，研究两类具有明确对比对象和验证方式的改进方法。创新点的表述重点强调“现有方法在哪一环节仍存在问题、本文具体改进哪一步、如何通过实验判断改进是否有效”。

\section{基于预测不确定性的轻量化风险感知运动基元规划}

动态障碍短时预测能够提高局部规划的前瞻性，但预测结果会受到感知误差、运动模型偏差以及障碍转向、加减速等行为变化影响。现有实时局部规划中，采用当前状态或单一未来轨迹进行碰撞判断的方法计算开销较小，但当预测逐渐失准时容易低估未来风险；机会约束、碰撞概率和风险优化等方法能够显式描述不确定性，但部分方法需要额外的概率计算或连续优化，直接用于高频局部规划时会增加计算负担。

针对这一矛盾，本研究拟在现有运动基元局部规划框架基础上，引入面向预测不确定性的低复杂度时空风险评价。利用动态障碍未来状态均值、协方差或等价安全包络，评价候选运动基元与未来障碍状态之间的风险关系，使预测不确定程度能够直接参与候选轨迹排序；同时保留TTC和最小安全距离等硬安全过滤，以避免仅依赖软风险代价。

因此，该部分的改进点不是提出新的通用预测模型，也不是重新设计运动基元，而是研究“不确定性信息如何以较低计算代价进入运动基元筛选”。后续通过无预测、确定性预测和加入不确定性风险评价后的运动基元规划进行对比，并在障碍变速、转向和预测误差增大条件下分析碰撞率、最小安全距离、任务成功率和规划耗时。若该方法能够在计算开销可接受的前提下，使性能在预测失准时退化更缓、碰撞风险更低，则说明这一改进具有实际价值。

\section{通信状态与邻机相关性联合驱动的协同约束自适应方法}

在单机局部规划基础上扩展到多无人机后，邻机状态和计划轨迹成为分布式冲突判断的重要信息来源。现有通信约束研究已经考虑连通性、通信对象选择、信息新鲜度、时延鲁棒性和通信质量自适应等问题，但通信指标与局部规划需求之间仍存在差异。通信质量较好并不意味着某个邻机一定会影响当前轨迹；反之，即使邻机信息存在一定时延，只要其与本机具有较强的潜在冲突趋势，该邻机仍可能需要重点处理。

针对这一问题，本研究拟将通信状态评价与邻机规划相关性结合起来。首先利用时延、丢包、信息年龄和链路状态描述共享信息的时效性与可靠程度；进一步结合机间距离、相对速度、潜在TTC和轨迹交叉趋势评价邻机对当前局部规划的相关程度，并据此筛选需要重点参与冲突判断的邻机。在此基础上，根据通信可靠程度和邻机相关性调整轨迹可信时间窗、机间安全裕度、状态外推范围及冲突约束强度；通信持续恶化时，再采用保守协同或局部自主作为安全退化方式。

因此，该部分的改进点不是提出新的通信指标，也不是单纯设计三级协同模式，而是研究“通信可靠程度+邻机规划相关程度”如何共同作用于协同对象筛选和规划约束调整。后续将与固定邻机集合/固定约束方法、仅通信质量驱动的方法进行对比，并通过去除邻机相关性筛选或去除协同约束调整的消融实验，分析任务完成率、机间最小距离、碰撞率、通信开销和规划时间。若该方法能够在通信恶化时减少无关邻机处理，并对关键冲突邻机保持更合理的安全约束，则说明这一改进能够改善通信受限条件下的安全性与效率折中。
